\chapter{Convex Optimisation}\label{ch:qm:convex-optimisation}

At a quarter to four, an optimiser must turn two hundred return forecasts into a long-short portfolio before the close: dollar-neutral and
neutral to each of five sectors, no position above 1.5\% of capital, gross exposure at most 100\%, and no more than 20\% of capital traded
against yesterday's book. It must return the best portfolio these limits allow, and it must also say what each limit costs, because tomorrow
someone will ask whether the turnover limit is worth it. On the day a risk manager adds a limit that contradicts the others, it must say that no
portfolio exists and show why, in a form a person can check. Convex optimisation does all three: a solution, a price for every constraint, and a
certificate when there is no solution. This chapter builds the theory that guarantees them (convexity, duality, the optimality conditions), the
problem classes that portfolio construction lives in, and a solver of the firm's own, and runs it on the portfolio.

\section{Convex sets, functions and problems}

\begin{definition}[Convex set, convex function, convex optimisation problem]\label{def:qm:convex-optimisation:convex}
A set $C \subseteq \R^n$ is a \emph{convex set}\index{convex set} if it contains the segment between any two of its points. A function $f$ is a \emph{convex
function}\index{convex function} if $f(\theta x + (1 - \theta)y) \le \theta f(x) + (1 - \theta)f(y)$ for $\theta \in [0, 1]$. A \emph{convex optimisation
problem}\index{convex optimisation problem} minimises a convex function subject to convex inequality constraints $g_i(x) \le 0$ and affine equality constraints
$Ax = b$.
\end{definition}

\begin{proposition}[Local minima are global]\label{prop:qm:convex-optimisation:local}
For a convex problem, every local minimum is a global minimum, and the set of minimisers is convex.
\end{proposition}

\begin{proof}
If $x$ is a local minimum and $y$ is feasible with $f(y) < f(x)$, the points $x + \theta(y - x)$ are feasible for $\theta \in (0, 1]$ and, by convexity, $f(x + \theta(y -
x)) \le f(x) + \theta(f(y) - f(x)) < f(x)$, however small $\theta$: a contradiction. If $x$ and $y$ both minimise, so does every point between them.
\end{proof}

This is why convexity matters more than any algorithm: a solver that stops at a point satisfying local conditions has found the answer. Portfolio problems are
convex when the objective is expected return minus a multiple of a variance built from a positive semidefinite covariance (chapter~22), and the constraints are
linear limits, absolute values bounded above (gross exposure, turnover) and norms bounded above (risk and tracking error).

\section{Duality}

\begin{definition}[Lagrangian, dual problem]\label{def:qm:convex-optimisation:dual}
For minimising $f(x)$ subject to $g_i(x) \le 0$ and $Ax = b$, the \emph{Lagrangian}\index{Lagrangian} is $L(x, \lambda, \nu) = f(x) + \sum_i\lambda_ig_i(x) + \nu^\top(Ax
- b)$, and the dual function is $d(\lambda, \nu) = \inf_xL(x, \lambda, \nu)$. The \emph{dual problem}\index{dual problem} maximises $d$ over $\lambda \ge 0$.
\end{definition}

\begin{proposition}[Weak duality]\label{prop:qm:convex-optimisation:weak}
For any feasible $x$ and any $\lambda \ge 0$, $d(\lambda, \nu) \le f(x)$; hence the dual optimum is at most the primal optimum.
\end{proposition}

\begin{proof}
For feasible $x$, $\lambda_ig_i(x) \le 0$ and $Ax - b = 0$, so $L(x, \lambda, \nu) \le f(x)$, and $d(\lambda, \nu) \le L(x, \lambda, \nu)$.
\end{proof}

\begin{definition}[Strong duality, Slater's condition, shadow price]\label{def:qm:convex-optimisation:strong}
\emph{Strong duality}\index{strong duality} holds when the primal and dual optimal values are equal. \emph{Slater's condition}\index{Slater's condition} is the existence of a
strictly feasible point ($g_i(x) < 0$ for the nonlinear constraints, $Ax = b$). The optimal dual variable of a constraint is its \emph{shadow price}\index{shadow price}: the
rate at which the optimal value changes when the constraint is relaxed.
\end{definition}

For a convex problem, \omterm{def:qm:convex-optimisation:strong}{Slater's condition} implies \omterm{def:qm:convex-optimisation:strong}{strong duality} and the existence of optimal dual variables. If the constraint $g_i(x) \le 0$ is loosened to $g_i(x) \le u$,
the optimal value $p^*(u)$ satisfies $\partial p^*/\partial u = -\lambda_i$ where it is differentiable: a multiplier is the marginal value of the limit, in the units of the
objective. Relaxation measures the same thing globally.

On the chapter's portfolio (\cref{fig:qm:convex-optimisation:turnover}), the objective is expected return minus five times the variance (annual units, covariance of
chapter~22). With a 20\% turnover limit the optimum has an expected return of 3.78\%, a volatility of 3.47\% and a utility of 3.18\%; the limit binds, and its multiplier is
0.0159: each extra 1\% of turnover allowed is worth 1.6 basis points of utility. Solving again with the limit at 10\%, 15\%, 30\% and 40\% traces the curve whose slope the
multiplier is: 0.0210 at 10\%, 0.0184 at 15\%, 0.0119 at 30\%, falling to 0.0011 at 80\% as the limit stops mattering.

\begin{omfigure}
\begin{tikzpicture}
\begin{axis}[width=10cm,height=5cm,xlabel={turnover limit (\% of capital)},ylabel={annual \%},
  tick label style={font=\small},label style={font=\small},xmin=0,xmax=85,ymin=2.6,ymax=4.4,grid=major,grid style={gray!25},
  legend columns=3,legend style={font=\footnotesize,at={(0.5,-0.3)},anchor=north,draw=none,fill=none,
  /tikz/every even column/.append style={column sep=6pt}},table/col sep=comma]
\addplot[omThm,very thick,mark=square*,mark size=1.5pt] table[x=limit,y=expret]{figdata/methods/23-convex-optimisation/turnover.csv};
\addplot[omDef,very thick,mark=*,mark size=1.5pt] table[x=limit,y=utility]{figdata/methods/23-convex-optimisation/turnover.csv};
\addplot[black,dashed,thick,domain=10:30,samples=2] {3.180 + 0.01593*(x-20)};
\legend{expected return,utility,slope at 20\% = multiplier}
\end{axis}
\end{tikzpicture}

\omcaption{The price of turnover: the optimal long-short portfolio's utility and expected return as the turnover limit is relaxed from 5\% to 80\% of capital. The dashed line
through the 20\% point has the slope of the turnover constraint's multiplier, 0.0159 (1.6 basis points of utility per 1\% of turnover). Data: the chapter's tutorial,
seeded.}\label{fig:qm:convex-optimisation:turnover}
\end{omfigure}

Relaxing each constraint group in turn and re-solving ranks them by what they cost (\cref{fig:qm:convex-optimisation:costs}): the turnover limit 0.44\% of utility, the gross
limit 0.28\%, the forty-three binding position limits together 0.19\%, sector neutrality 0.014\%, dollar neutrality nothing. The last is instructive: dollar neutrality is
implied by the five sector neutralities, which sum to it, so it is redundant, and its multiplier (0.0004 in the solver's answer) is not unique; a zero relaxation value with a
nonzero multiplier is the signature of a redundant constraint.

\begin{omfigure}
\begin{tikzpicture}
\begin{axis}[width=9cm,height=4.8cm,xbar,bar width=9pt,
  xlabel={utility lost to the constraint (annual \%)},
  ytick={0,1,2,3,4},yticklabels={dollar neutrality,sector neutrality,position limits,gross exposure,turnover},
  tick label style={font=\small},label style={font=\small},
  xmin=0,xmax=0.55,ymin=-0.6,ymax=4.6,grid=major,grid style={gray!25},
  xticklabel style={/pgf/number format/fixed,/pgf/number format/precision=2},scaled x ticks=false,table/col sep=comma]
\addplot[fill=omDef!60,draw=omDef,nodes near coords,nodes near coords style={font=\footnotesize,anchor=west,xshift=1pt,
  /pgf/number format/fixed,/pgf/number format/precision=3}]
  table[x=relax,y=k]{figdata/methods/23-convex-optimisation/costs.csv};
\end{axis}
\end{tikzpicture}

\omcaption{What each constraint group costs the long-short portfolio: the gain in utility when it is removed and the problem re-solved. Dollar neutrality costs nothing because the
sector neutralities imply it. Data: the chapter's tutorial, seeded.}\label{fig:qm:convex-optimisation:costs}
\end{omfigure}

\begin{definition}[Infeasibility certificate]\label{def:qm:convex-optimisation:certificate}
An \emph{infeasibility certificate}\index{infeasibility certificate} for the linear system $Ax = b$, $Gx \le h$ is a pair $(\lambda \ge 0, \nu)$ with $G^\top\lambda + A^\top\nu = 0$ and
$h^\top\lambda + b^\top\nu < 0$.
\end{definition}

By Farkas' lemma a certificate exists if and only if the system is infeasible: combining the rows with the weights $\lambda$ and $\nu$ gives $0 = (\lambda^\top G + \nu^\top A)x \le
h^\top\lambda + b^\top\nu < 0$, an impossibility anyone can verify. A phase-one problem, minimise $t$ subject to $Ax = b$, $Gx \le h + t\mathbf 1$, produces one from its dual variables when
its optimum is positive. When a risk manager demands that each of the forty stocks of the first sector be held at least 0.5\% long, while sector neutrality is on, the phase-one
optimum is $t^* = 0.005$ and the certificate puts weight $1/40$ on each of the forty minimum-position rows and a multiplier on the first sector's neutrality equation, and nothing on any
other limit: adding the forty rows gives ``the sector is at least 20\% long'', and the neutrality row says it is exactly 0. The report names the two rules that clash.

\section{Optimality conditions}

\begin{definition}[Karush--Kuhn--Tucker conditions]\label{def:qm:convex-optimisation:kkt}
The \emph{Karush--Kuhn--Tucker conditions}\index{Karush--Kuhn--Tucker conditions} at $(x, \lambda, \nu)$ are primal feasibility, dual feasibility ($\lambda \ge 0$), complementary
slackness ($\lambda_ig_i(x) = 0$) and stationarity, $\nabla f(x) + \sum_i\lambda_i\nabla g_i(x) + A^\top\nu = 0$ (Karush, 1939; Kuhn and Tucker, 1951).
\end{definition}

\begin{theorem}[KKT conditions characterise the optimum]\label{thm:qm:convex-optimisation:kkt}
For a convex problem with differentiable functions satisfying \omterm{def:qm:convex-optimisation:strong}{Slater's condition}, $x$ is optimal if and only if there are $(\lambda, \nu)$ such that $(x, \lambda, \nu)$ satisfies
the KKT conditions; the $(\lambda, \nu)$ are then optimal dual variables.
\end{theorem}

\begin{proof}
Sufficiency: stationarity says $x$ minimises the \omterm{def:qm:convex-optimisation:convex}{convex function} $L(\cdot, \lambda, \nu)$, so $d(\lambda, \nu) = L(x, \lambda, \nu) = f(x)$ by complementary slackness and
feasibility; weak duality then makes $x$ primal optimal and $(\lambda, \nu)$ dual optimal. Necessity: under Slater, \omterm{def:qm:convex-optimisation:strong}{strong duality} holds with an optimal $(\lambda, \nu)$, and
$f(x) = d(\lambda, \nu) \le L(x, \lambda, \nu) \le f(x)$ forces complementary slackness and the minimisation of $L$ at $x$, whose first-order condition is stationarity.
\end{proof}

Complementary slackness is the economics: a limit that does not bind has a zero price, and a limit with a positive price binds. The solver's answer for the portfolio satisfies the
conditions to working precision: stationarity residual and complementarity both below $10^{-9}$. The forty-three positions at the 1.5\% limit carry positive multipliers,
about 0.0175 each; the other 157 carry zero.

\section{Quadratic and second-order cone programmes}

\begin{definition}[Linear, quadratic, second-order cone and semidefinite programmes]\label{def:qm:convex-optimisation:classes}
A \emph{linear programme}\index{linear programme} minimises a linear function subject to linear constraints; a \emph{quadratic programme}\index{quadratic programme} has a convex quadratic
objective and linear constraints; a \emph{second-order cone programme}\index{second-order cone programme} (SOCP) adds constraints $\lVert A_ix + b_i\rVert_2 \le c_i^\top x + d_i$; a
\emph{semidefinite programme}\index{semidefinite programme} constrains a symmetric matrix depending affinely on $x$ to be positive semidefinite.
\end{definition}

Each class contains the previous one. Gross exposure and turnover, sums of absolute values, become linear by splitting each position into a positive and a negative part: the chapter's
portfolio of 200 stocks is a \omterm{def:qm:convex-optimisation:classes}{quadratic programme} in 1\,000 variables.

\begin{proposition}[Risk and impact as cone constraints]\label{prop:qm:convex-optimisation:cones}
A volatility or tracking-error limit $\sqrt{(w - b)^\top\Sigma(w - b)} \le \sigma_{\max}$ is the second-order cone constraint $\lVert L^\top(w - b)\rVert_2 \le \sigma_{\max}$ with $\Sigma =
LL^\top$. A market-impact cost growing like the 3/2 power of the traded amount, $|x|^{3/2} \le s$, is representable with two rotated cone constraints: it holds if and only if there is
$r \ge 0$ with $x^2 \le rs$ and $r^2 \le |x|$.
\end{proposition}

\begin{proof}
The first is the identity $w^\top\Sigma w = \lVert L^\top w\rVert^2$. For the second (with $x \ge 0$ after splitting), if $x^{3/2} \le s$ take $r = \sqrt x$: then $r^2 = x$ and $x^2 = \sqrt x\,x^{3/2} \le rs$.
Conversely $r^2 \le x$ gives $r \le \sqrt x$, so $x^2 \le rs \le \sqrt x\,s$ and $x^{3/2} \le s$. Each constraint $u^2 \le vw$ with $v, w \ge 0$ is a rotated second-order cone,
$\lVert(2u, v - w)\rVert \le v + w$.
\end{proof}

A risk limit and a risk penalty are the same problem seen from two sides. Maximising expected return subject to a 4\% volatility limit (dollar-neutral, positions within 1.5\%) by the
firm's cone solver, and maximising return minus $\gamma/2$ times variance with $\gamma$ found by bisection to give 4\% volatility, produce the same portfolio (the weights differ by
$3 \times 10^{-8}$), an expected return of 5.31\%, $\gamma = 25.8$, and a cone multiplier of 1.033 $= \gamma \times 0.04$: the \omterm{def:qm:convex-optimisation:strong}{shadow price} of the volatility limit, expected return per unit of
volatility allowed.

\begin{definition}[Nearest correlation matrix]\label{def:qm:convex-optimisation:ncm}
The \emph{nearest correlation matrix}\index{nearest correlation matrix} to a symmetric matrix $C$ is the positive semidefinite matrix with unit diagonal closest to $C$ in the Frobenius norm, a
\omterm{def:qm:convex-optimisation:classes}{semidefinite programme} that Higham (2002) solves by alternating projections with Dykstra's correction.
\end{definition}

Correlations estimated pair by pair from data with gaps are not guaranteed to form a correlation matrix, and a covariance built from them can make an optimiser promise negative variance.
Fifty series of sixty days with half the observations missing give a pairwise matrix with 21 negative eigenvalues, the smallest $-1.42$; the \omterm{def:qm:convex-optimisation:ncm}{nearest correlation matrix}, reached in 74
projections, is 4.45 away in the Frobenius norm, changes no entry by more than 0.31, and has no negative eigenvalue (\cref{fig:qm:convex-optimisation:ncm}).

\begin{omfigure}
\begin{tikzpicture}
\begin{axis}[width=10cm,height=4.8cm,xlabel={eigenvalue rank},ylabel={eigenvalue},
  tick label style={font=\small},label style={font=\small},xmin=0,xmax=51,ymin=-1.6,ymax=3,grid=major,grid style={gray!25},
  legend style={font=\footnotesize,at={(0.03,0.97)},anchor=north west,fill=white,draw=none},table/col sep=comma]
\addplot[omThm,thick,mark=*,mark size=1.2pt] table[x=rank,y=before]{figdata/methods/23-convex-optimisation/ncm.csv};
\addplot[omDef,thick,mark=square*,mark size=1.2pt] table[x=rank,y=after]{figdata/methods/23-convex-optimisation/ncm.csv};
\addplot[black,thin,sharp plot] coordinates {(0,0) (51,0)};
\legend{pairwise correlations,nearest correlation matrix}
\end{axis}
\end{tikzpicture}

\omcaption{Eigenvalues (in increasing order; the largest, above 10, off the chart) of a pairwise-complete correlation matrix of fifty series with half the data missing, and of its \omterm{def:qm:convex-optimisation:ncm}{nearest
correlation matrix}. Twenty-one negative eigenvalues become zero. Data: the chapter's tutorial, seeded.}\label{fig:qm:convex-optimisation:ncm}
\end{omfigure}

\section{What solvers actually do}

\begin{definition}[Interior-point method]\label{def:qm:convex-optimisation:ipm}
An \emph{interior-point method}\index{interior-point method} follows the central path: it applies Newton's method to the KKT conditions with complementary slackness relaxed to $s_iz_i = \mu$
for slacks $s$ and multipliers $z$, keeping both strictly positive, and drives $\mu$ to zero.
\end{definition}

Each iteration solves one linear system in the Newton direction; Mehrotra's (1992) predictor-corrector computes an affine direction, uses it to choose how far to reduce $\mu$, and corrects
for the second-order term, with the same matrix factorisation. The number of iterations barely depends on the problem's size: the chapter's portfolio converges in 18, the residuals
and the complementarity falling from order one to below $10^{-9}$ (\cref{fig:qm:convex-optimisation:ipm}). \omterm{def:qm:convex-optimisation:ipm}{Interior-point methods} descend from Karmarkar's (1984) polynomial algorithm for
\omterm{def:qm:convex-optimisation:classes}{linear programmes} and were extended to cone programmes by Nesterov and Nemirovskii (1994). First-order splitting methods such as ADMM, of which the OSQP solver (Stellato and coauthors, 2020)
is a production example, need only one factorisation and cheap iterations, reach moderate accuracy fast, and warm-start well when yesterday's solution is close to today's; the firm's
cone solver is of this kind, and needed 924 iterations for the 4\% volatility limit above.

\begin{omfigure}
\begin{tikzpicture}
\begin{semilogyaxis}[width=10cm,height=5cm,xlabel={interior-point iteration},ylabel={residual},
  tick label style={font=\small},label style={font=\small},xmin=0,xmax=18,ymin=1e-10,ymax=100,grid=major,grid style={gray!25},
  legend columns=3,legend style={font=\footnotesize,at={(0.5,-0.3)},anchor=north,draw=none,fill=none,
  /tikz/every even column/.append style={column sep=8pt}},table/col sep=comma]
\addplot[omDef,very thick,mark=*,mark size=1.2pt] table[x=iter,y=dual]{figdata/methods/23-convex-optimisation/ipm.csv};
\addplot[omThm,very thick,mark=square*,mark size=1.2pt] table[x=iter,y=primal]{figdata/methods/23-convex-optimisation/ipm.csv};
\addplot[black,thick,dashed] table[x=iter,y=mu]{figdata/methods/23-convex-optimisation/ipm.csv};
\legend{stationarity,primal feasibility,complementarity $\mu$}
\end{semilogyaxis}
\end{tikzpicture}

\omcaption{Convergence of the firm's interior-point solver on the long-short portfolio (1\,000 variables): the norms of the stationarity and feasibility residuals and the average complementarity
$\mu$ at each Mehrotra iteration, floored at $10^{-9}$ (below it, the last digits depend on the order of floating-point operations). Data: the chapter's tutorial, seeded.}\label{fig:qm:convex-optimisation:ipm}
\end{omfigure}

\section{Tutorial: the constraint that cost the most}\label{tut:qm:convex-optimisation:shadow}

\begin{tutorial}
\textbf{Goal.} Formulate the long-short portfolio as a \omterm{def:qm:convex-optimisation:classes}{quadratic programme}, solve it with the firm's \omterm{def:qm:convex-optimisation:ipm}{interior-point method}, verify the KKT conditions, read the \omterm{def:qm:convex-optimisation:strong}{shadow prices}, and read the
certificate of an infeasible set of limits.
\textbf{End state:} \cref{fig:qm:convex-optimisation:turnover,fig:qm:convex-optimisation:costs,fig:qm:convex-optimisation:ipm}; the turnover multiplier 0.0159 and the ranking of the
constraints.
\begin{tutsteps}
\item \textbf{The Newton step} of the \omterm{def:qm:convex-optimisation:ipm}{interior-point method}, with the predictor and the corrector sharing one matrix.
\omcode{code/firm/portopt/firm_portopt.py}{50}{81}{The interior-point iteration: Mehrotra's predictor and corrector share one matrix.}\label{lst:qm:convex-optimisation:ipm}
\item \textbf{The phase-one problem} and its Farkas certificate.
\omcode{code/firm/portopt/firm_portopt.py}{89}{107}{Phase one and the infeasibility certificate.}\label{lst:qm:convex-optimisation:farkas}
\item \textbf{Run} \texttt{today()}, \texttt{constraint\_costs()}, \texttt{turnover\_curve()}, \texttt{infeasible()}, \texttt{risk\_limit()} and \texttt{broken\_correlation()} in
\texttt{qm\_opt.py}, then \texttt{fig\_opt.py}.
\end{tutsteps}
\textbf{What to change next.} Replace the turnover limit by a 3/2-power impact cost in the objective and compare the traded amounts; warm-start the \omterm{def:qm:convex-optimisation:ipm}{interior-point method} from yesterday's solution
and count the iterations saved.
\end{tutorial}

\section{Build: the portfolio optimiser}\label{bld:qm:convex-optimisation:portopt}

\begin{build}
\bfield{Purpose} The miniature firm's portfolio construction: forecasts, a covariance (chapter~22) and the limits in, orders out, with a price for each limit and a certificate when the limits
contradict each other.
\bfield{Interface} \texttt{qp(P, q, A, b, G, h)}; \texttt{farkas(A, b, G, h)}; \texttt{admm(P, q, C, lower, upper, cones)}; \texttt{PortfolioProblem(alpha, Sigma, gamma, w0)} with
\texttt{add\_budget}, \texttt{add\_neutral}, \texttt{add\_box}, \texttt{add\_gross}, \texttt{add\_turnover} and \texttt{solve()}; \texttt{nearest\_correlation(C)}.
\bfield{Rules} Every solve reports its status, KKT residuals and the dual variable of each named constraint; an infeasible request returns the certificate, never a silently relaxed portfolio;
covariances pass a positive-semidefiniteness check (and the nearest-correlation repair) before use.
\bfield{Acceptance tests} \texttt{code/firm/portopt/tests/}: a small QP and its KKT residuals; the minimum-variance closed form and its multiplier; Farkas certificates on feasible and infeasible
systems; the cone solver on a norm ball; portfolio constraints hold and the turnover multiplier equals the finite-difference slope; the \omterm{def:qm:convex-optimisation:ncm}{nearest correlation matrix} is a correlation matrix and
leaves one unchanged.
\bfield{Stretch} A homogeneous self-dual embedding (certificates without phase one); sparse KKT factorisation for thousands of assets; the 3/2-power impact cost in the cone solver.
\end{build}

\begin{omsources}
\item W.~Karush, \emph{Minima of Functions of Several Variables with Inequalities as Side Constraints}, MSc thesis, University of Chicago, 1939; H.~W.~Kuhn and A.~W.~Tucker, ``Nonlinear
programming'', \emph{Second Berkeley Symposium}, 1951.
\item H.~Markowitz, ``Portfolio selection'', \emph{Journal of Finance} 7, 1952.
\item N.~Karmarkar, ``A new polynomial-time algorithm for linear programming'', \emph{Combinatorica} 4, 1984; S.~Mehrotra, ``On the implementation of a primal-dual interior point method'',
\emph{SIAM Journal on Optimization} 2, 1992.
\item Y.~Nesterov and A.~Nemirovskii, \emph{Interior-Point Polynomial Algorithms in Convex Programming}, SIAM, 1994.
\item S.~Boyd and L.~Vandenberghe, \emph{Convex Optimization}, Cambridge University Press, 2004.
\item M.~S.~Lobo, L.~Vandenberghe, S.~Boyd and H.~Lebret, ``Applications of second-order cone programming'', \emph{Linear Algebra and its Applications} 284, 1998.
\item N.~J.~Higham, ``Computing the nearest correlation matrix: a problem from finance'', \emph{IMA Journal of Numerical Analysis} 22, 2002.
\item B.~Stellato, G.~Banjac, P.~Goulart, A.~Bemporad and S.~Boyd, ``OSQP: an operator splitting solver for quadratic programs'', \emph{Mathematical Programming Computation} 12, 2020.
\end{omsources}

\section{Exercises}

\begin{exercise}[$\star$]\label{exo:qm:convex-optimisation:1}
Is $f(w) = \sqrt{w^\top\Sigma w}$ convex for positive semidefinite $\Sigma$? Is $\{w: w^\top\Sigma w \ge 0.01\}$ a \omterm{def:qm:convex-optimisation:convex}{convex set}?
\end{exercise}

\begin{exercise}[$\star$]\label{exo:qm:convex-optimisation:2}
The turnover multiplier is 0.0159. Estimate the utility gained by raising the limit from 20\% to 22\%, and compare with re-solving (the curve gives 3.180\% at 20\% and a slope of 0.0139 over
[20\%, 30\%]).
\end{exercise}

\begin{exercise}[$\star$]\label{exo:qm:convex-optimisation:3}
Write $\sum_i|w_i - w_{0,i}| \le \tau$ with linear constraints only.
\end{exercise}

\begin{exercise}[$\star\star$]\label{exo:qm:convex-optimisation:4}
Derive the \omterm{def:qm:covariance-estimation-and-random-matrices:pca}{minimum-variance portfolio} and its multiplier from the KKT conditions of $\min\frac12w^\top\Sigma w$ subject to $\mathbf 1^\top w = 1$.
\end{exercise}

\begin{exercise}[$\star\star$]\label{exo:qm:convex-optimisation:5}
Find a Farkas certificate for $x_1 + x_2 = 1$, $x_1 \le 0.2$, $x_2 \le 0.3$.
\end{exercise}

\begin{exercise}[$\star\star$]\label{exo:qm:convex-optimisation:6}
Show that $\max\alpha^\top w$ subject to $w^\top\Sigma w \le \sigma^2$ (no other constraint) has solution $w = \sigma\Sigma^{-1}\alpha/\sqrt{\alpha^\top\Sigma^{-1}\alpha}$ and a \omterm{def:qm:convex-optimisation:strong}{shadow price} of the
volatility limit equal to $\sqrt{\alpha^\top\Sigma^{-1}\alpha}$.
\end{exercise}

\begin{exercise}[$\star\star\star$]\label{exo:qm:convex-optimisation:7}
\emph{Coding.} Rank the constraint groups of the chapter's portfolio by the utility each costs, by relaxing them one at a time, and explain the result for dollar neutrality.
\end{exercise}

\begin{exercise}[$\star\star\star$]\label{exo:qm:convex-optimisation:8}
\emph{Find the flaw.} ``The optimiser failed tonight, so the overnight job dropped constraints one at a time until it found a feasible portfolio, and traded that.''
\end{exercise}

\section{Problem: The Constraint That Cost the Most}

\begin{problem}[{Weekend problem --- pricing the limits of a long-short book}]\label{pb:qm:convex-optimisation:1}
A book of 200 stocks (the factor covariance of chapter~22, annual units) maximises expected return minus five times the variance, dollar- and sector-neutral, positions within 1.5\%, gross
exposure at most 100\%, turnover at most 20\% against yesterday's optimal book. Today's forecasts are yesterday's decayed by 0.8 plus news.

\medskip\noindent\textbf{Part I --- The optimum.}
\begin{enumerate}
\item What are the optimal portfolio's expected return, volatility and utility?
\item Which constraints bind, and how many positions sit at their limits?
\item How many interior-point iterations does the solver need, and what are the KKT residuals?
\item What is the turnover limit's multiplier, and what does it mean?
\item What is the gross-exposure multiplier?
\end{enumerate}

\medskip\noindent\textbf{Part II --- Prices.}
\begin{enumerate}[resume]
\item How does the utility change as the turnover limit goes from 5\% to 80\%?
\item How do the multipliers compare with the slopes of that curve?
\item Rank the constraint groups by the utility each costs.
\item Why is the dollar-neutrality cost zero although its multiplier is not?
\item What is the expected return given up per 1\% of turnover?
\end{enumerate}

\medskip\noindent\textbf{Part III --- When there is no portfolio.}
\begin{enumerate}[resume]
\item What happens when sector 0's stocks must each be at least 0.5\% long?
\item What does the certificate say, and how is it checked?
\item How would a volatility limit of 4\% be expressed, and what is its \omterm{def:qm:convex-optimisation:strong}{shadow price}?
\item Why must the correlation matrix be checked before optimising?
\item What does the \omterm{def:qm:convex-optimisation:ncm}{nearest correlation matrix} do to a pairwise matrix with 21 negative eigenvalues?
\end{enumerate}

\medskip\noindent\textbf{Part IV --- Judgement.}
\begin{enumerate}[resume]
\item Which limit would you ask the risk committee to loosen first?
\item Is a 1.6-basis-point price per 1\% of turnover worth paying, and against what?
\item What should the optimiser's report contain every day?
\item State the \emph{named result}: the \omterm{def:qm:convex-optimisation:strong}{shadow price} of the turnover limit and the ranking of the constraints by the return they cost.
\item In one sentence: what does duality give a portfolio manager that the solution alone does not?
\end{enumerate}
\end{problem}

\section{Interview questions}

\begin{interviewq}[$\star$ \iqroles{researcher, developer}]\label{iq:qm:convex-optimisation:1}
Why is convexity so important in optimisation?
\end{interviewq}

\begin{interviewq}[$\star\star$ \iqroles{researcher, trader}]\label{iq:qm:convex-optimisation:2}
What is a \omterm{def:qm:convex-optimisation:strong}{shadow price}, and how would you use one to argue for loosening a risk limit?
\end{interviewq}

\begin{interviewq}[$\star\star$ \iqroles{researcher}]\label{iq:qm:convex-optimisation:3}
State the KKT conditions and explain complementary slackness in portfolio terms.
\end{interviewq}

\begin{interviewq}[$\star\star$ \iqroles{developer, researcher}]\label{iq:qm:convex-optimisation:4}
Your optimiser says a problem is infeasible. How do you find out which constraints conflict?
\end{interviewq}

\begin{interviewq}[$\star\star$ \iqroles{developer}]\label{iq:qm:convex-optimisation:5}
Interior-point or first-order methods for a daily portfolio rebalance of 3\,000 stocks: which, and why?
\end{interviewq}

\begin{interviewq}[$\star\star\star$ \iqroles{researcher}]\label{iq:qm:convex-optimisation:6}
How do you write a tracking-error limit and a 3/2-power impact cost as convex constraints?
\end{interviewq}
